\chapter*{Glossaire}
\addcontentsline{toc}{section}{Glossaire}

\begin{description}
    \item[Système multi-agents (SMA)] Approche informatique consistant à modéliser un système comme un ensemble d'entités autonomes (agents) qui interagissent dans un environnement partagé selon leurs propres objectifs et stratégies.
    
    \item[Marché boursier] Place de marché organisée où s'échangent des valeurs mobilières, notamment des actions d'entreprises cotées.
    
    \item[Trading algorithmique] Utilisation d'algorithmes informatiques pour automatiser les décisions d'achat et de vente sur les marchés financiers.
    
    \item[Agent] Entité logicielle autonome capable de percevoir son environnement, de prendre des décisions et d'agir pour atteindre ses objectifs.
    
    \item[Ontologie] Représentation formelle et partagée des connaissances d'un domaine, définissant les concepts, leurs propriétés et leurs relations.
    
    \item[PASSI] Process for Agent Societies Specification and Implementation, méthodologie de conception structurée pour les systèmes multi-agents.
    
    \item[UML] Unified Modeling Language, langage de modélisation graphique standardisé utilisé en génie logiciel.
    
    \item[Protocole d'interaction] Ensemble de règles définissant la manière dont les agents communiquent et coordonnent leurs actions.
    
    \item[Dynamiques émergentes] Phénomènes complexes apparaissant au niveau global d'un système résultant des interactions entre ses composants individuels.
    
    \item[Courtier] Intermédiaire financier facilitant les transactions entre acheteurs et vendeurs sur un marché.
    
    \item[Régulateur] Agent ou institution chargé de surveiller et contrôler les opérations du marché pour garantir son bon fonctionnement et sa conformité aux règles établies.
\end{description}